\scanpage{292}
% archive sheet number: 293
et que ces suites et la suite constante de valeur 1 forment une base de l\deleted{\unsure{a}}'espace vectoriel $E_r$ de toutes les suites complexes de période $r$. Montrer que la série
\[
u(s)_n c_n
\]
est alternée à valeurs absolues décroissantes tendant vers zéro, donc convergente. En déduire une autre démonstration de b).

d) Soit $z$ un nombre complexe, et posons $u_n=\unsure{z|Z}^{n}$ pour $n\quad 1$. \deleted{\unsure{D}}Montrer que la suite des sommes partielles $S_n$ de (1) reste bornée si et seulement si on a
\[
z\quad1\ ,\quad z\ne1\tag*{(5)}
\]
(auquel cas la suite des $S_n$ est majorée par $\;2\,\dfrac{z}{1-z}\;$). \deleted{Montrer que si la série $c_n$ est divergente, alors} \inserted{En déduire que} la condition (5) implique que la série \deleted{$\unsure{u_n c_n|c_n}$} \inserted{$c_n z^n$} est convergente. Montrer qu'elle est même absolument convergente pour $z\quad1$, tandis que pour un $z$ tel que $\;z\;=1$, $z\ne1$, on a convergence absolue si et seulement si la \deleted{suite} série $c_n$ est convergente. Enfin, \deleted{\unsure{Application} à la convergence et à la convergence absolue\unsure{:} montrer que pour toute partie $A$} \inserted{$A$ désignant une partie} de l'ensemble \inserted{$E$} des nombres complexes satisfaisant (5), pour que la convergence de la série $c_n z^n$ soit uniforme pour $z\quad A$, il suffit qu'il existe un nombre $a\quad0$ tel que $z\quad A$ implique $1-z\quad a$. Montrer que \deleted{\unsure{c}}cette condition est aussi nécessaire si la série $c_n$ est divergente (tandis que si $c_n$ est convergente, la convergence est uniforme pour l'ensemble des nombres complexes $z$ tels que $\;z\quad1$). \deleted{\unsure{E}}

e) Montrer que la fonction $f(z)$ sur l'ensemble $E$ des nombres complexes $z$ satisfaisant (5), représentée par la série
\[
f(z)=c_n z^n\tag*{(6)}
\]
est continue. Elle admet une limite finie pour $z=1$ si et seulement si la \deleted{suite} série $c_n$ est convergente, auquel cas la fonction $f(z)$ définie par la formule (6) est définie et continue \deleted{pour} sur tout le disque unité $z\quad1$.

f) Appliquer les résultats de d) et e) à la série trigonométrique
\[
c_n e^{ix}\quad ,
\]
et aux séries
\[
c_n\cos nx\ ,\quad c_n\sin nx\quad :
\]
en particulier, ces séries sont convergentes pour tout $x$ réel n'appartenant pas à l'ensemble \deleted{$\unsure{2\ \mathrm R|2\ \mathrm Z}$} $2\quad\underline{\mathrm Z}$ (i.e. $x$ pas de la forme $2\ n$, avec $n$ entier), et représentent des fonctions continues de $x$ pour $x\quad\underline{\mathrm R}-2\quad\underline{\mathrm Z}$ .

g) Montrer que la suite $(z^n)_{n\quad1}$ ($z$ étant un nombre complexe) est périodique de période $r$ si et seulement si $z$ est une racine r.ème de l'unité. \deleted{\unsure{Lorsque}} Déduire de ceci et et b) que lorsque $z$ est une racine de l'unité la série \deleted{\unsure{En}} $c\ z^n$ est convergente, \unsure{dauf|sauf} dans le cas où $z=1$ et $c_n$ \unsure{est}
